La première inversion conceptuelle consiste à concevoir le vide comme une réponse relationnelle positive.\par

Dans HMT, le vide est la réponse dynamique minimale du système lorsque ses deux feuillets doivent coexister, se propager et clore leur différence en conservant l’orientation qui les distingue.\par

Les deux feuillets constituent deux lectures conjuguées du même antécédent angulaire. L’un conserve une orientation et l’autre conserve l’orientation opposée. C’est pourquoi apparaissent deux canaux, \(q_+\) et \(q_-\), comme les deux faces spectrales corrélées d’un même état.\par

L'opérateur\par

\[
\mathcal V_{90,120}
=
(I-T^{90})(I-T^{120})^{-1}
\]

compare deux profondeurs de propagation : \(90\) et \(120\). Ces profondeurs sont déterminées par le même mode élémentaire :\par

\[
90=3\cdot30,
\qquad
120=4\cdot30.
\]

C’est pourquoi la réponse de chaque feuillet peut se réécrire comme\par

\[
\frac{1-q^{90}}{1-q^{120}}
=
\frac{1+s+s^2}{1+s+s^2+s^3},
\qquad
s=q^{30}.
\]

Et c’est ici que l’interprétation devient réellement puissante. Le numérateur contient trois états du mode \(30\) ; le dénominateur contient ces mêmes trois états et un de plus. Le vide est donc une complétion exacte \(3\to4\).\par

Ce quatrième terme intègre et conserve les trois précédents. Il les complète. Ce qui est ajouté est enregistré comme un résidu positif de complétion :\par

\[
1-\frac{1+s+s^2}{1+s+s^2+s^3}
=
\frac{s^3}{1+s+s^2+s^3}>0.
\]

La mémoire du vide est précisément l’accroissement exact qui permet à trois termes de devenir quatre. Elle constitue l’empreinte positive et nécessaire de la complétion.\par

Des deux valeurs propres de cet unique opérateur découlent les quatre relations constitutives :\par

\[
\widehat\varepsilon=r_+^2,
\qquad
\widehat\mu=r_-^2,
\qquad
\widehat Z=\frac{r_-}{r_+},
\qquad
\widehat c=\frac1{r_+r_-}.
\]

Cela change complètement l’interprétation de la perméabilité et de la permittivité.\par

La permittivité est la réponse quadratique d’un feuillet. La perméabilité est la réponse quadratique du feuillet conjugué. L’impédance conserve la relation orientée entre les deux. La vitesse de propagation dépend de leur produit commun.\par

Ainsi, produit et quotient accomplissent deux tâches différentes :\par

\begin{itemize}[leftmargin=*,itemsep=0.35em]
\item le produit oublie quel était chaque feuillet et conserve la propagation commune ;
\item le quotient conserve quelle était leur orientation relative.
\end{itemize}

Le produit répond à la question : « que partagent les deux feuillets ? ». Le quotient répond à la question : « comment se distinguent-ils ? ».\par

Por eso\par

\[
\widehat\mu\,\widehat\varepsilon=\widehat c^{-2}
\]

y\par

\[
\frac{\widehat\mu}{\widehat\varepsilon}
=
\widehat Z^2
\]

sont les deux manières complémentaires de lire la même paire spectrale.\par

L’aspect holographique apparaît ici de façon très nette : quatre grandeurs visibles ne contiennent en réalité que deux degrés de liberté internes. Chacune de certaines paires adéquates permet de reconstruire les deux feuillets. La publication constitutive conserve la généalogie lorsque produit et quotient restent conjointement disponibles.\par

Telle est la réparation profonde de l’ablation de \(\mu\) et de \(\varepsilon\). La réparation réunit leurs valeurs et leur origine opératorielle commune.\par

La carte du Système international peut ensuite transformer ces sections intrinsèques en \(\mu_0\), \(\varepsilon_0\), \(Z_0\) et \(c\). Sa fonction consiste à exprimer, au sein d’une convention métrologique, une structure dont la provenance a déjà été sélectionnée par l’opérateur.\par

Le vide cesse ainsi d’être un réservoir de quatre constantes. Il devient une unique opération spectrale avec deux mémoires conjuguées.\par

Et cette vision a une conséquence philosophique importante : le vide est une relation de compatibilité entre deux manières orientées de réaliser un même état.\par

